% Compile from the FSI repository root:
%   bash docs/resources/benchmarks/Model2_bluntbased_report/build_report.sh
% Or:
%   pdflatex -output-directory docs/resources/benchmarks/Model2_bluntbased_report/build \
%     docs/resources/benchmarks/Model2_bluntbased_report/model2_bluntbased_current_report.tex
\documentclass[11pt,a4paper]{article}

\usepackage[utf8]{inputenc}
\usepackage[T1]{fontenc}
\usepackage[italian]{babel}
\usepackage{geometry}
\usepackage{graphicx}
\usepackage{subcaption}
\usepackage{booktabs}
\usepackage{siunitx}
\usepackage{amsmath}
\usepackage{amssymb}
\usepackage{float}
\usepackage{xcolor}
\usepackage{hyperref}
\usepackage{pgfplots}
\pgfplotsset{compat=1.18}
\pgfplotsset{every axis/.append style={grid=both, grid style={gray!20}, tick align=outside}}

% Resolve asset paths from repo root OR from this report folder.
\makeatletter
\newcommand{\fsiroot}{}
\newcommand{\fsipath}[1]{\fsiroot#1}
\newcommand{\reportpath}[1]{figures/#1}
\IfFileExists{figures/mode_shapes_alpha0.png}{%
  % compile from report folder
}{%
  \IfFileExists{docs/resources/benchmarks/Model2_bluntbased_report/figures/mode_shapes_alpha0.png}{%
    \renewcommand{\reportpath}[1]{docs/resources/benchmarks/Model2_bluntbased_report/figures/#1}%
  }{%
    \IfFileExists{../../../../docs/resources/benchmarks/Model2_bluntbased_report/figures/mode_shapes_alpha0.png}{%
      \renewcommand{\reportpath}[1]{../../../../docs/resources/benchmarks/Model2_bluntbased_report/figures/#1}%
    }{}%
  }%
}
\IfFileExists{output_data/Model2-bluntbased/Model2-bluntbased_dry_egv_eigendata.csv}{%
  \def\fsiroot{}%
}{%
  \IfFileExists{../../../../output_data/Model2-bluntbased/Model2-bluntbased_dry_egv_eigendata.csv}{%
    \def\fsiroot{../../../../}%
  }{%
    \PackageWarning{model2_bluntbased}{Cannot locate FSI repo root; figures/CSV may be missing.}%
  }%
}
\makeatother

\geometry{margin=24mm}
\sisetup{detect-all=true, per-mode=symbol}

\title{Stato attuale della replica numerica del Modello~2 (strut blunt-based)}
\author{Report tecnico di lavoro --- \texttt{Model2-bluntbased}}
\date{27 maggio 2026 --- allineato a \texttt{config\_Model2-bluntbased.py} e log \texttt{20260527\_105737}}

\begin{document}
\maketitle

\begin{abstract}
Questo report documenta lo stato corrente della replica del \textbf{Modello~2} (strut con
trailing edge smussato) del benchmark sperimentale di Bash e Liu,
\emph{Bending Flutter and Torsional Flutter of Flexible Hydrofoil Struts}.
Documenta gli \textbf{input attivi} in \texttt{config\_Model2-bluntbased.py}, l'analisi
modale a secco, la massa aggiunta strip a fluido fermo, l'accoppiamento PanelAero
(\texttt{nspan=20}, \texttt{nchord=10}, quartic) e un primo sweep Roger RFA
(\texttt{V\_list}, 50 velocita'). I risultati numerici citati provengono dal log
\texttt{flutter\_analysis\_Model2-bluntbased\_20260527\_105737.log}, salvo il riepilogo
flutter da \texttt{output\_data/flutter\_summary.csv} (run precedente completata).
Confronto con curve Bash \& Liu: da fare.
\end{abstract}

\section{Riferimento sperimentale}

Il documento di riferimento descrive quattro modelli di strutture flessibili
(Modello~1--4) con diverso rapporto di massa e geometria di corda. Il \textbf{Modello~2}
e' la variante con \emph{blunt trailing edge} (sezione A-A del disegno di costruzione:
circonferenze al leading edge, facce piatte al trailing edge, $t/c \approx 1/30$).
Nel setup sperimentale la strut e' immersa, con sweep di \SI{15}{\degree} e campata
nominalmente di \SI{49.963}{in} (\SI{1.269}{\meter}) e corda \SI{11.59}{in}
(\SI{0.294}{\meter}).

In questa fase del lavoro numerico si replica \textbf{soltanto la strut} (nessun piano
di portanza aggiuntivo): beam cantilever con sezione definita in
\texttt{config\_Model2-bluntbased.py}, accoppiata a una superficie DLM equivalente
costruita in \texttt{FLUID/PanelAero/Qjj/executer.py}.

\section{Configurazione numerica}

La configurazione attiva e' in
\texttt{examples/Model2-bluntbased/config\_Model2-bluntbased.py}.
I default globali (\texttt{examples/config.py}) impongono \texttt{n\_elements=200},
\texttt{num\_modes\_egv=4}, fluido acqua, \texttt{nspan=20}, \texttt{nchord=10}.

\subsection{Geometria e materiali}

\begin{table}[H]
  \centering
  \caption{Parametri principali della configurazione \texttt{Model2-bluntbased}.}
  \label{tab:config}
  \begin{tabular}{lll}
    \toprule
    Grandezza & Valore config & Nota \\
    \midrule
    Campata $L$ & $49.963\times0.0254 = \SI{1.269}{\meter}$ & \texttt{beam\_length}, non proiettata \\
    Corda $c$ & $5.796\times2\times0.0254 = \SI{0.294}{\meter}$ & \texttt{chord}, 12\,inches nominali \\
    $t/c$ & $1/30$ & \texttt{thickness\_factor=0.033} \\
    Sweep $\Lambda$ & \SI{15}{\degree} & \texttt{sweep\_angle\_deg} \\
    Tip mesh (SC) & $(x,y,z)\approx(\SI{0.340}{\meter},\SI{1.269}{\meter},0)$ & $x=L\tan\Lambda$; \texttt{...\_nodes.csv} \\
    $\mu$ (lineare) & \SI{18.8}{\kilogram\per\meter} & \texttt{mu} \\
    $EI_{xx}$, $GJ$ & \SI{3710.91}, \SI{5940.9}{\newton\meter\squared} & flessione / torsione \\
    $EA$, $EI_{zz}$, $GA_x$, $GA_z$ & $\times 1000$ rispetto a $EI_{xx}$/$GJ$ & rigidezze assiali/scherro moltiplicate \\
    $I_\alpha$ & \SI{0.16812}{lbf.in.s$^2$} & MOI totale su EA (strut + tip) \\
    $k_{i22}$, $k_{i33}$ & 0.3, 0.7 & inerzie rotazionali secondarie \\
    $x_{SC}/c$, $x_{CM}/c$ & 0.31, 0.417 & \texttt{xea\_factor}, \texttt{xcm\_factor} \\
    Offset $e$ (sezione $M$) & \SI{0.0304}{\meter} & log: da SC a CG lungo corda \\
    Dihedral & \SI{-90}{\degree} & solo trave (\texttt{dihedral\_rotate\_aerogrid=False}) \\
    $\zeta$ Rayleigh & 0.5\% (2 modi) & \texttt{rayleigh\_target\_zetas} \\
    Elementi FEA & 200 & default \texttt{examples/config.py} \\
    Modi EGV / flutter & 4 / 4 & \texttt{num\_modes\_egv}, \texttt{num\_modes\_flutter\_egv} \\
    \bottomrule
  \end{tabular}
\end{table}

\subsection{Geometria della trave nel codice}

I nodi sono su \texttt{beam\_reference\_node\_positions} (\texttt{beam\_geometry.py}):
campata lungo \textbf{+Y}, sweep $x(y)=y\tan\Lambda$, radice SC in $(0,0,0)$.
L'ultima mesh esportata ha $y_\mathrm{tip}=L$ e $x_\mathrm{tip}\approx L\tan15^\circ$
(coerente con \texttt{beam\_length} fisico nel config).
\texttt{xea\_factor}/\texttt{xcm\_factor} definiscono la sezione $6\times6$ e il
coupling PanelAero (\texttt{xea\_factor=0.31} sui corner della griglia), non lo
spostamento dei nodi trave.
Nel modello custom (\texttt{structural\_model.py}) gli offset in $M$ usano
$c/\cos\Lambda$; il log riporta $x_{SC}=\SI{0.088}{\meter}$, $x_{CM}=\SI{0.119}{\meter}$,
$e=\SI{0.0304}{\meter}$.
La strip theory usa $c$ dal config e passo $\Delta y=L/n_\mathrm{el}$.

\subsection{Orientamento del modello}

Dopo la rotazione di dihedral \SI{-90}{\degree} sulla sola trave, la campata operativa
nel sistema ruotato e' lungo l'asse \texttt{beam\_span\_axis=2} (asse $Z$ del riferimento
post-rotazione). L'analisi modale e le forme modali sono quindi visualizzate con
ascissa \emph{lunghezza d'arco lungo la trave} (coerente con la Figura delle mode
shapes a $\alpha=0^\circ$).

\section{Griglia aerodinamica (PanelAero / DLM)}

L'analisi accoppiata carica la griglia pre-calcolata indicata nel config:
\begin{center}
\small
\texttt{qjj\_precomputed/Model2-bluntbased\_water\_alpha0.0\_nspan20\_nchord10\_quartic/aerogrid.npz}
\end{center}
(200 pannelli, metodo \texttt{quartic}, $\alpha=0^\circ$, acqua).
Parametri allineati al config: \texttt{nspan=20}, \texttt{nchord=10},
\texttt{xea\_factor=0.31} per l'asse elastico sui corner DLM.

Il blocco \texttt{executer.py} per \texttt{Model2-bluntbased} (pre-calcolo \texttt{Qjj}
manuale) usa ancora geometria \emph{proiettata} e risoluzione \texttt{22$\times$11}:
campata $L\cos15^\circ$, corda $c/\cos15^\circ$. Va aggiornato per coerenza con il config
prima di rigenerare le \texttt{Qjj}.

\begin{figure}[H]
  \centering
  \includegraphics[width=\linewidth]{\reportpath{aerogrid_beam_alpha0.png}}
  \caption{Griglia aerodinamica DLM sovrapposta alla trave (centro di taglio e baricentro).
  Punti di forza al \SI{25}{\percent} della corda, punti di controllo al \SI{75}{\percent}
  (convenzione PanelAero). Vista generata da \texttt{executer.py} con
  \texttt{plot\_aero\_beam\_model}.}
  \label{fig:aerobeam}
\end{figure}

Statistiche accoppiamento (log 27/05): distanza media spline trave--pannello
\SI{0.090}{\meter}, max \SI{0.202}{\meter}; 200 CP / 200 punti di forza.

\section{Analisi modale a secco}

\subsection{Modello FEA e vincoli}

\begin{itemize}
  \item Trave a \texttt{n\_elements = 200} elementi;
  \item Primi \texttt{num\_modes\_egv = 4} modi elastici;
  \item Nodo radice: tutti e 6 i DOF vincolati (\texttt{Constrained DOFs: [0..5]});
  \item Matrici ridotte: $M,K$ di dimensione $(1200\times1200)$ dopo eliminazione vincoli.
\end{itemize}

I risultati completi sono in
\texttt{output\_data/Model2-bluntbased/}:
\texttt{Model2-bluntbased\_dry\_egv\_eigendata.csv},
\texttt{...\_nodes.csv}, \texttt{...\_constrained\_dofs.csv}.

\subsection{Frequenze proprie a secco}

\begin{table}[H]
  \centering
  \caption{Frequenze proprie a secco (analisi statica / eigenvalue).}
  \label{tab:dry-freq}
  \begin{tabular}{ccc}
    \toprule
    Modo & $f_\mathrm{dry}$ [Hz] & Commento qualitativo \\
    \midrule
    1 & 4.586 & prevalenza $w$ e $\phi_x$ \\
    2 & 27.936 & flessione/torsione accoppiata \\
    3 & 40.486 & contributo marcato $\phi_y$ \\
    4 & 75.200 & modo superiore, forma oscillante in $w$ \\
    \bottomrule
  \end{tabular}
\end{table}

\subsection{Forme modali}

\begin{figure}[H]
  \centering
  \includegraphics[width=\linewidth]{\reportpath{mode_shapes_alpha0.png}}
  \caption{Forme modali normalizzate ($\alpha = 0^\circ$): spostamenti $u,v,w$ e rotazioni
  $\phi_x,\phi_y,\phi_z$ lungo la campata. I DOF $u$, $v$ e $\phi_z$ sono trascurabili
  ($\sim 10^{-9}$--$10^{-10}$); la dinamica e' dominata da $w$, $\phi_x$ e $\phi_y$.}
  \label{fig:modes}
\end{figure}

\section{Massa aggiunta strip e frequenze a fluido fermo}

Dopo la riduzione modale, il codice aggiunge la massa aggiunta strip (cilindro
equivalente, $\rho = \SI{997}{\kilogram\per\meter\cubed}$) sui DOF trasversali
e risolve il sistema ridotto con smorzamento (state-space a $V=0$).

Con corda fisica $c\approx\SI{0.294}{\meter}$, la strip theory fornisce
$m_a'=\rho\pi c^2/4 \approx \SI{67.88}{\kilogram\per\meter}$
(vs.\ $\mu=\SI{18.8}{\kilogram\per\meter}$ strutturale; integrale lump
$m_a' L \approx \SI{86}{\kilogram}$ su $L\approx\SI{1.269}{\meter}$).
La versione precedente con corda proiettata $c/\cos15^\circ$ dava erroneamente
$m_a'\approx\SI{72.5}{\kilogram\per\meter}$.

Dal log \texttt{flutter\_analysis\_Model2-bluntbased\_20260527\_105737.log}:
\begin{itemize}
  \item $m_a' = \SI{67.8843}{\kilogram\per\meter}$,
    passo lump $\Delta y = L/n_\mathrm{el} \approx \SI{6.35e-3}{\meter}$,
    $\approx \SI{0.215}{\kilogram}$ per DOF trasversale per nodo;
  \item \texttt{strip\_theory\_span\_axis=1};
  \item \texttt{dihedral\_deg=-90} impostato ma \textbf{non} applicato a $M_\mathrm{added}$
    (proiezione modale con autovettori dry nel telaio di riferimento);
  \item $\|M_\mathrm{add}\|_F \approx 4.31$, $\| \hat{M}_\mathrm{added}\|_F \approx 3.00$.
\end{itemize}

\subsection{Frequenze a fluido fermo}

\begin{table}[H]
  \centering
  \caption{Confronto frequenze dry vs.\ sistema a fluido fermo con strip theory.}
  \label{tab:wet-freq}
  \begin{tabular}{cccc}
    \toprule
    Modo & $f_\mathrm{dry}$ [Hz] & $f_\mathrm{fluid\ rest}$ [Hz] & $\Delta f / f_\mathrm{dry}$ \\
    \midrule
    1 & 4.586 & 2.767 & $-40\%$ \\
    2 & 27.936 & 17.156 & $-39\%$ \\
    3 & 40.486 & 38.717 & $-4\%$ \\
    4 & 75.200 & 47.484 & $-37\%$ \\
    \bottomrule
  \end{tabular}
\end{table}

\begin{table}[H]
  \centering
  \caption{Massa aggiunta strip generalizzata diagonale $\hat{M}_{\mathrm{added},ii}$.}
  \label{tab:madd}
  \begin{tabular}{cc}
    \toprule
    Modo & $\hat{M}_{\mathrm{added},ii}$ [kg] \\
    \midrule
    1 & 1.746 \\
    2 & 1.631 \\
    3 & 0.276 \\
    4 & 1.509 \\
    \bottomrule
  \end{tabular}
\end{table}

I valori sono salvati in
\texttt{model2\_bluntbased\_frequencies.csv} e
\texttt{model2\_bluntbased\_strip\_added\_mass.csv} in questa cartella.

\begin{figure}[H]
  \centering
  \begin{tikzpicture}
    \begin{axis}[
      ybar,
      bar width=14pt,
      width=0.88\linewidth,
      height=0.52\linewidth,
      xlabel={Numero modo},
      ylabel={Frequenza [Hz]},
      ymin=0, ymax=82,
      symbolic x coords={1,2,3,4},
      xtick=data,
      enlarge x limits=0.18,
      legend style={font=\small, at={(0.98,0.98)}, anchor=north east},
      grid=both,
      grid style={gray!20},
      tick align=outside
    ]
      \addplot[fill=blue!55, draw=blue!80!black] coordinates {
        (1,4.586) (2,27.936) (3,40.486) (4,75.200)
      };
      \addlegendentry{Dry}
      \addplot[fill=red!55, draw=red!80!black] coordinates {
        (1,2.767) (2,17.156) (3,38.717) (4,47.484)
      };
      \addlegendentry{Fluido fermo (strip)}
    \end{axis}
  \end{tikzpicture}
  \caption{Confronto a barre delle frequenze proprie dry e a fluido fermo (strip theory)
  per ciascun modo.}
\end{figure}

\section{Primo sweep flutter (Roger RFA)}

Config: \texttt{V\_list = linspace(0.5, 45, 50)} m/s, \texttt{n\_lag=4},
\texttt{blag = [0.5, 1.0, 1.5, 2.0]} (adimensionali, convertiti in $1$/m nel solver),
\texttt{empirical\_fluid\_damping=True}, modello \texttt{abramson\_delta} con
$\Delta\delta = \SI{5}{\per\second}$, \texttt{aero\_im\_Q\_scale=1.0}.

Riepilogo da \texttt{output\_data/flutter\_summary.csv} (analisi completata, 26 maggio;
da rivalutare dopo allineamento definitivo DLM--struttura):

\begin{table}[H]
  \centering
  \caption{Velocita' di flutter stimate per branch modale (primo attraversamento $g=0$).}
  \label{tab:flutter}
  \begin{tabular}{cccc}
    \toprule
    Branch & $V_f$ [m/s] & $V_f$ [kn] & $f_f$ [Hz] \\
    \midrule
    1 & --- & --- & --- \\
    2 & 14.86 & 28.88 & 41.24 \\
    3 & 14.61 & 28.41 & 39.17 \\
    4 & 10.56 & 20.52 & 67.06 \\
    \bottomrule
  \end{tabular}
\end{table}

Il branch~4 ha la velocita' critica minima ($\approx\SI{10.6}{\meter\per\second}$).
Il branch~1 non mostra attraversamento nel riepilogo.

\begin{figure}[H]
  \centering
  \includegraphics[width=\linewidth]{\reportpath{vg_vf_combined.png}}
  \caption{Diagrammi V--$g$ e V--$f$ dal run con \texttt{flutter\_summary.csv} (26 maggio).}
  \label{fig:vgvf}
\end{figure}

\section{Pipeline idro-aeroelastica (stato lavori)}

\begin{table}[H]
  \centering
  \caption{Stato dei componenti per il benchmark Model2.}
  \begin{tabular}{lll}
    \toprule
    Componente & Stato & File / nota \\
    \midrule
    Config struttura + sezione blunt & OK & \texttt{config\_Model2-bluntbased.py} \\
    EGV dry + export CSV & OK & \texttt{output\_data/Model2-bluntbased/} \\
    Mode shapes / aerobeam plot & OK & Fig.\ \ref{fig:modes}, \ref{fig:aerobeam} \\
    Strip $M_\mathrm{added}$ + freq.\ fluid rest & OK & Tab.\ \ref{tab:wet-freq}, log 27/05 \\
    \texttt{Qjj} usate in analisi & OK & \texttt{nspan20\_nchord10\_quartic} \\
    \texttt{executer.py} CAERO1 Model2 & Da allineare & ancora $L\cos\Lambda$, \texttt{22$\times$11} \\
    Mesh Capytaine & Presente & \texttt{FLUID/capytaine/Model2-bluntbased/} \\
    Sweep flutter Roger RFA & Parziale & Tab.\ \ref{tab:flutter}; rigenerare con geom.\ finale \\
    Confronto Bash \& Liu & Da fare & digitazione curve \\
    \bottomrule
  \end{tabular}
\end{table}

Altri parametri dal config:
\begin{itemize}
  \item Capytaine: \texttt{omega\_list} 0.1--150 rad/s (300 punti),
    mesh \texttt{30$\times$30}, \texttt{results\_modal\_radiation/};
  \item sezione blunt: contorno normalizzato \texttt{raw} (41 punti, TE piatto);
  \item plot: \texttt{dimensionless\_vgvf\_results}, contributi $K$/$C$, partecipanti DLM.
\end{itemize}

\section{Conclusione e prossimi passi}

Config e analisi strutturale sono allineati su $L=\SI{1.269}{\meter}$,
$c=\SI{0.294}{\meter}$, $\mu=\SI{18.8}{\kilogram\per\meter}$,
$m_a'=\SI{67.88}{\kilogram\per\meter}$. Le \texttt{Qjj} in uso
(\texttt{20$\times$10}) seguono il config; \texttt{executer.py} per rigenerarle no.
Un primo flutter numerico indica $V_f\approx\SI{10.6}{\meter\per\second}$ (branch~4),
da ricalcolare dopo allineamento DLM.

\textbf{Prossimi passi:}
\begin{enumerate}
  \item Aggiornare \texttt{executer.py} (geom.\ fisica, \texttt{20$\times$10}) e rigenerare \texttt{Qjj}.
  \item Completare sweep $V$--$g$--$f$ con la geometria finale e confrontare Bash \& Liu.
  \item Digitare curve sperimentali in CSV (stile report Abramson).
  \item Opzionale: applicare dihedral a $M_\mathrm{added}$ strip.
\end{enumerate}

\end{document}
